% !TeX program = xelatex
\documentclass[12pt,a4paper]{article}
\usepackage{fontspec}
\setmainfont{Times New Roman}[Path=C:/Windows/Fonts/,UprightFont=times.ttf,BoldFont=timesbd.ttf,ItalicFont=timesi.ttf,BoldItalicFont=timesbi.ttf]
\usepackage{amsmath,amssymb,amsfonts}
\usepackage[left=1.5cm,right=1.5cm,top=1.5cm,bottom=1.7cm,headheight=17pt]{geometry}
\usepackage{enumitem,array,graphicx,tikz,fancyhdr,needspace}
\usetikzlibrary{arrows.meta,calc,patterns}
% Sửa số đề và số trang đề thực tế sau khi biên dịch/tách PDF.
\newcommand{\soDe}{1}
\newcommand{\soTrangDe}{1}
\newcommand{\maDe}{\ifnum\soDe<1000 0\fi\ifnum\soDe<100 0\fi\ifnum\soDe<10 0\fi\soDe}
\pagestyle{fancy}
\fancyhf{}
\lhead{\small\textbf{HƯỚNG TỚI THPTQG 2027}}
\rhead{\small\textit{Biên soạn: Dương Gia Bảo}}
\rfoot{\small Trang \thepage/\soTrangDe\ - Mã đề: \maDe}
\renewcommand{\headrulewidth}{0.4pt}
\fancypagestyle{answers}{\rfoot{\small Trang \thepage\ - Mã đề: \maDe}}
\fancypagestyle{firstpage}{\fancyhead{}\renewcommand{\headrulewidth}{0pt}}
\setlength{\parindent}{0pt}
\setlength{\parskip}{3pt}
\setlength{\emergencystretch}{1em}
\setlist[enumerate,1]{label=\textbf{Câu \arabic*:},leftmargin=*,itemsep=8pt,topsep=5pt}
\setlist[enumerate,2]{label=\alph*),leftmargin=*,itemsep=3pt,topsep=3pt}
\newcommand{\bonpa}[4]{\par\nobreak\smallskip\nobreak\noindent
\begin{tabular}{@{}p{0.47\linewidth}p{0.47\linewidth}@{}}
\textbf{A.} #1 & \textbf{B.} #2\\[8pt]
\textbf{C.} #3 & \textbf{D.} #4\end{tabular}}
\newcommand{\loigiai}[1]{\par\medskip\Needspace{4\baselineskip}\noindent\textbf{#1}\par\nobreak}
\begin{document}
\thispagestyle{firstpage}
\noindent
\begin{minipage}[t]{0.35\textwidth}
    \centering
    {\fontsize{13}{15}\selectfont\textbf{HƯỚNG TỚI THPTQG 2027}}\\[5pt]
    \textit{Biên soạn: Dương Gia Bảo}\\[6pt]
    \setlength{\fboxrule}{0.8pt}%
    \setlength{\fboxsep}{5pt}%
    \fbox{\begin{minipage}{0.83\linewidth}
        \centering
        {\fontsize{13}{15}\selectfont\textbf{ĐỀ ÔN TẬP}\\[2pt]\textbf{CHÍNH THỨC}}\\[5pt]
        \textit{(Đề thi gồm có \ifnum\soTrangDe<10 0\fi\soTrangDe\ trang)}
    \end{minipage}}
\end{minipage}\hfill
\begin{minipage}[t]{0.63\textwidth}
    \centering
    {\fontsize{13}{15}\selectfont\textbf{KỲ KHẢO SÁT CHẤT LƯỢNG LỚP 12}\\[2pt]
    \textbf{NĂM HỌC 2026 - 2027}}\\[6pt]
    {\fontsize{14}{16}\selectfont\textbf{Môn: TOÁN - Lớp 12 (ĐỀ SỐ \soDe)}}\\[5pt]
    {\fontsize{11.5}{13}\selectfont\textit{Thời gian làm bài: 90 phút, không kể thời gian giao đề}}\\[6pt]
    \setlength{\fboxrule}{0.8pt}%
    \setlength{\fboxsep}{6pt}%
    \fbox{\makebox[3cm]{\fontsize{14}{16}\selectfont\textbf{Mã đề: \maDe}}}
\end{minipage}
\par\vspace{3pt}
\hrule height 0.6pt
\vspace{4pt}
\noindent\textit{Họ và tên học sinh:}\dotfill\hspace{1cm}\textit{Lớp:}\makebox[3.5cm]{\dotfill}
\par\vspace{3pt}
\hrule height 0.6pt
\vspace{4pt}

% Chèn 12 câu Phần I (dùng \bonpa), 4 câu Phần II, 6 câu Phần III và dòng HẾT.
% Nếu kèm đáp án: sau \newpage, đặt \pagestyle{answers} trước BẢNG ĐÁP ÁN TỔNG HỢP.
% HƯỚNG DẪN GIẢI CHI TIẾT cũng bắt đầu ở trang mới.
\end{document}
